% -----------------------------------------------------------------------------
% Active stationary-hull closure: no active degenerate stationary-family
% profile is Seregin-attainable
% -----------------------------------------------------------------------------
% This section replaces the remaining active stationary-hull closure obligation by an explicit
% consequence of the local terminal concentration decomposition/no-splitting
% package.  It uses no global coercive estimate on the original Navier--Stokes
% flow.  The
% only inputs are local compactness, local terminal profile decomposition,
% inactive/radiative/rough discharge, finite active no-splitting, mildness
% inheritance, and the Albritton--Barker endpoint Liouville theorem.
% -----------------------------------------------------------------------------

\section{Exclusion of active degenerate stationary-hull profiles by terminal stratification}
\label{sec:R3-terminal-stratification-closure}

We close the remaining analytic statement in the degenerate stationary-hull
class by the terminal local concentration theorem proved in the main
paper.  The point is that the obstruction is not specific to the existence of a
non-isolated stationary family.  Once a stationary-family element is attained as
an active Type~I Seregin limit, it is an active stationary terminal carrier.
Such carriers are excluded by the same local terminal concentration analysis and
no-splitting mechanism used for the final residual class.  No global coercive
estimate for the original solution is used.

Throughout this section, a centered Type~I profile satisfies
\begin{equation}\label{eq:R3S-centered}
   \partial_\tau U-\Delta U+\frac12 y\cdot\nabla U+\frac12 U
   +(U\cdot\nabla)U+\nabla \Pi=0,
   \qquad \nabla\cdot U=0
\end{equation}
on \(\mathbb R^3\times\mathbb R\), with the usual local suitability and
pressure-gauge conventions.

\subsection{The terminal local package}

We use the following specialization of the terminal local concentration theorem
proved in the main text.  It is stated here exactly in the form needed for the
active stationary-hull closure.

\begin{theorem}[Terminal local exclusion package]
\label{thm:R3S-terminal-package}
Let \((U,\Pi)\) be a bounded centered Type~I Seregin profile satisfying local
suitability and carrying a retained local CKN concentration, i.e. for some
compact parabolic cylinder \(Q\) and some \(\varepsilon_*>0\),
\begin{equation}\label{eq:R3S-retained-activity}
   \iint_Q\bigl(|U|^3+|\Pi-c_Q(\tau)|^{3/2}\bigr)\,dy\,d\tau
   \ge \varepsilon_* .
\end{equation}
Assume that \((U,\Pi)\) is outside the previously closed small-amplitude,
stationary \(L^3\), tight, fast-decay, structured, axisymmetric
bounded-circulation, and rotational relative-equilibrium classes.  Then the
following hold.
\begin{enumerate}[label=(\roman*)]
   \item Every terminal extraction of \((U,\Pi)\) belongs to one of the standard
   local alternatives: a single retained profile, scattering, exterior,
   radiative, rough, or a finite separated family of retained active profiles.
   \item Scattering, exterior, radiative, and rough terminal alternatives cannot
   be the sole carriers of the retained compact CKN concentration
   \eqref{eq:R3S-retained-activity}.
   \item Finite separated families of retained active profiles are impossible.
   \item Any remaining singleton active terminal carrier is terminally atomic:
   it has no retained active tail descendant and no radiative or rough
   descendant.
   \item Every retained active terminally atomic carrier has a backward sequence
   of finite critical norm: there exist \(\tau_k\to-\infty\) such that
   \begin{equation}\label{eq:R3S-atomic-L3-input}
      \sup_k \|U(\cdot,\tau_k)\|_{L^3(\mathbb R^3)}<\infty .
   \end{equation}
   \item Retained active terminal profiles inherit mildness under the physical
   pullback.  Therefore the Albritton--Barker endpoint Liouville theorem applies
   to any profile satisfying \eqref{eq:R3S-atomic-L3-input}.
\end{enumerate}
\end{theorem}

\begin{proof}
This is \Cref{thm:terminal-local-stratification-package} applied to
\((U,\Pi)\).  The local alternatives in (i) are those of
\Cref{thm:terminal-exhaustion-main}.  Item (ii) is the non-active stratum
discharge \Cref{thm:terminal-nonactive-discharge}.  Item (iii) is
\Cref{thm:no-finite-separated-profile-family}.  Item (iv) is
\Cref{prop:no-multi-implies-atomic}.  Item (v) is
\Cref{lem:atomic-sequence-L3}, and item (vi) is
\Cref{hyp:mildness-inheritance-main,thm:AB-main}.  These results are proved in
the present paper from Paper 0, Paper I, and local concentration analysis;
no global bound is assumed.
\end{proof}

\begin{remark}[Local nature of the package]
\label{rem:R3S-local-nature}
The theorem above is a local concentration statement.  It uses CKN compactness,
pressure-normalized local convergence, terminal profile decomposition, local
exclusion of inactive/radiative/rough carriers, and finite active
no-splitting.  It does not assert or require a global a priori bound for the
original physical solution.
\end{remark}

For completeness, and to make the dependence on local estimates explicit, we
record the elementary part of the package which converts terminal atomicity into
sequence-\(L^3\) control.

\begin{lemma}[Atomic local carriers have backward sequence-\(L^3\) control]
\label{lem:R3S-atomic-L3}
Let \((U,\Pi)\) be a bounded smooth ancient solution of
\eqref{eq:R3S-centered}.  Suppose that \((U,\Pi)\) is retained active and
terminally atomic: it has no retained active tail descendant, no radiative
descendant, and no rough descendant.  Then there exists a sequence
\(\tau_k\to-\infty\) such that
\begin{equation}\label{eq:R3S-atomic-L3}
   \sup_k \|U(\cdot,\tau_k)\|_{L^3(\mathbb R^3)}<\infty .
\end{equation}
\end{lemma}

\begin{proof}
We argue by contradiction.  Assume that no sequence satisfying
\eqref{eq:R3S-atomic-L3} exists.  Then
\begin{equation}\label{eq:R3S-L3-divergence}
   \liminf_{\tau\to-\infty}\|U(\cdot,\tau)\|_{L^3(\mathbb R^3)}=\infty .
\end{equation}
Choose \(\tau_n\to-\infty\) such that
\(
   \int_{\mathbb R^3}|U(y,\tau_n)|^3\,dy\to\infty .
\)
Since \(U\) is bounded, we may choose radii \(R_n\to\infty\) so that
\begin{equation}\label{eq:R3S-tail-mass}
   \int_{|y|>R_n}|U(y,\tau_n)|^3\,dy\ge 1 .
\end{equation}
Let
\begin{equation}\label{eq:R3S-unit-tail-sup}
   m_n:=\sup_{|x|>R_n}\int_{B_1(x)}|U(y,\tau_n)|^3\,dy .
\end{equation}
If \(\limsup_n m_n\ge\varepsilon_*\), then for a subsequence there exist
\(|x_n|>R_n\) such that
\begin{equation}\label{eq:R3S-active-tail-cube}
   \int_{B_1(x_n)}|U(y,\tau_n)|^3\,dy\ge \varepsilon_* .
\end{equation}
After recentering at \((x_n,\tau_n)\), either the local compactness and pressure
control required for a terminal profile fail, which is a rough descendant, or a
subsequence converges in the terminal topology to a retained active tail
descendant.  Both alternatives contradict terminal atomicity.

It remains to consider the case \(m_n<\varepsilon_*\) for all sufficiently large
\(n\).  Then \eqref{eq:R3S-tail-mass} says that critical mass persists outside
expanding balls, while \eqref{eq:R3S-unit-tail-sup} says that every unit-scale
exterior recentering remains below the retained-active threshold.  This is precisely the
radiative descendant alternative, again contradicting terminal atomicity.
Thus \eqref{eq:R3S-L3-divergence} is impossible, and
\eqref{eq:R3S-atomic-L3} follows.
\end{proof}

\subsection{Active stationary carriers are impossible}

\begin{theorem}[No active stationary residual carrier is Seregin-attainable]
\label{thm:R3S-no-active-stationary-carrier}
Let \((W,\Pi)\) be a stationary solution of the centered equation,
\begin{equation}\label{eq:R3S-stationary-centered}
   -\Delta W+\frac12 y\cdot\nabla W+\frac12 W
   +(W\cdot\nabla)W+\nabla\Pi=0,
   \qquad \nabla\cdot W=0 .
\end{equation}
Assume that:
\begin{enumerate}[label=(\roman*)]
   \item \((W,\Pi)\) is obtained as a centered Type~I Seregin limit of a
   finite-energy suitable weak solution at a singular point;
   \item \((W,\Pi)\) is retained active in the sense of
   \eqref{eq:R3S-retained-activity};
   \item \((W,\Pi)\) is outside all previously closed lower classes listed in
   \Cref{thm:R3S-terminal-package};
\end{enumerate}
Then no such \((W,\Pi)\) exists.
\end{theorem}

\begin{proof}
Suppose that such a stationary profile \((W,\Pi)\) exists.  Since it is a
centered Type~I Seregin limit and is retained active, it is an admissible input
for the terminal local package.  Apply \Cref{thm:R3S-terminal-package} to
\((W,\Pi)\).

Finite active terminal splittings are impossible by the terminal no-splitting
part of the package.  Scattering, exterior, radiative, and rough alternatives
cannot carry the retained compact CKN concentration.  Hence the retained
activity must be carried by a singleton terminal carrier.  By the atomicity
conclusion in \Cref{thm:R3S-terminal-package}, this singleton carrier is
terminally atomic.

Because \((W,\Pi)\) is stationary in \(\tau\), the singleton active carrier may
be represented by \((W,\Pi)\) itself after the terminal normalization.  Therefore
\Cref{lem:R3S-atomic-L3} gives a sequence \(\tau_k\to-\infty\) such that
\begin{equation}\label{eq:R3S-W-L3-sequence}
   \sup_k\|W\|_{L^3(\mathbb R^3)}
   =
   \sup_k\|W(\cdot,\tau_k)\|_{L^3(\mathbb R^3)}
   <\infty .
\end{equation}
Thus \(W\in L^3(\mathbb R^3)\).

Let
\begin{equation}\label{eq:R3S-physical-pullback}
   w(x,t)=(-t)^{-1/2}W\!\left(\frac{x}{\sqrt{-t}}\right),
   \qquad
   q(x,t)=(-t)^{-1}\Pi\!\left(\frac{x}{\sqrt{-t}}\right),
   \qquad t<0 .
\end{equation}
By mildness inheritance in the terminal package, \((w,q)\) is a mild ancient
Navier--Stokes solution after any finite time shift.  The critical scaling gives
\begin{equation}\label{eq:R3S-critical-L3-scaling}
   \|w(\cdot,t)\|_{L^3(\mathbb R^3)}=\|W\|_{L^3(\mathbb R^3)}<\infty
   \qquad (t<0).
\end{equation}
Applying the Albritton--Barker endpoint Liouville theorem along any sequence
\(t_k\downarrow-\infty\), we obtain \(w\equiv0\), and hence \(W\equiv0\).
This contradicts retained activity, since the pressure-normalized local CKN
quantity of the zero profile is zero.  Therefore no active stationary residual
carrier can be Seregin-attainable.
\end{proof}

\subsection{Closure of the degenerate stationary-hull residual class}

\begin{corollary}[No active degenerate stationary-family profile is Seregin-attainable]
\label{cor:R3S-no-active-degenerate-family}
Fix an admissible stationary phase space \(\mathfrak B\).  There is no stationary
centered profile \((W,\Pi)\) satisfying all of the following conditions:
\begin{enumerate}[label=(\roman*)]
   \item \(W\in\mathfrak B\), \(W\) solves the stationary centered equation
   \eqref{eq:R3S-stationary-centered}, and \(W\) lies on an admissible
   nontrivial stationary family;
   \item \((W,\Pi)\) is a centered Type~I Seregin limit of a finite-energy
   suitable weak solution at a singular point;
   \item \((W,\Pi)\) is retained active;
   \item \((W,\Pi)\) is outside the previously closed lower classes;
\end{enumerate}
\end{corollary}

\begin{proof}
The stationary-family condition in (i) is additional structure, but the exclusion in
\Cref{thm:R3S-no-active-stationary-carrier} applies to every active stationary
residual carrier.  Hence it applies in particular to stationary profiles lying
in admissible nontrivial stationary families.
\end{proof}

\begin{theorem}[Closure of the active degenerate stationary-hull class]
\label{thm:R3S-R3-empty}
Let \(\Cstat(\mathcal S;I,J;\mathfrak B)\) be the active
degenerate stationary-hull residual class.  Then
\begin{equation}\label{eq:R3S-R3-empty}
   \Cstat(\mathcal S;I,J;\mathfrak B)=\varnothing .
\end{equation}
\end{theorem}

\begin{proof}
Suppose, for contradiction, that
\(V\in\Cstat(\mathcal S;I,J;\mathfrak B)\).  By the
definition of the active degenerate stationary-hull class, the centered
translation hull of \((V,P)\) contains an active stationary profile \((W,\Pi)\)
which lies in an admissible nontrivial stationary family and is outside the
previously closed lower classes.  The scale-reselection lemma for centered hull
limits shows that \((W,\Pi)\) is itself a centered Type~I Seregin limit of the
original finite-energy suitable weak solution at the same singular point.

Thus \((W,\Pi)\) satisfies the hypotheses of
\Cref{cor:R3S-no-active-degenerate-family}.  This is impossible.  Hence no such
\(V\) exists, and \eqref{eq:R3S-R3-empty} follows.
\end{proof}

\begin{remark}[How this should be recorded in the residual decomposition]
\label{rem:R3S-decomposition}
The third residual line should be recorded as a theorem closed by the terminal
local concentration analysis:
\[
   \text{active degenerate stationary-hull residual}
   \quad\Longrightarrow\quad
   \text{active stationary residual carrier}
   \quad\Longrightarrow\quad
   \varnothing .
\]
The second implication is proved by terminal local concentration analysis:
non-atomic carriers produce an excluded active splitting, radiative descendant,
or rough descendant; atomic carriers have backward sequence-\(L^3\) control and
are ruled out by the endpoint ancient Liouville theorem.
\end{remark}
